The matter of the translation and rotation of the spherical harmonic functions is treated in its entirety by \textsc{Gumerov} \& \textsc{Duraiswami} in a plethora of articles on the topic of multipole expansions and their swift evaluation by the fast multipole method.
Chief among them is the paper wherein these rotations and translations are derived in detail,%
\footnote{%
  \cite{gumerov2003recursions} \textsc{Gumerov} \& \textsc{Duraiswami} (2003)
}
their book on the fast multipole method applied to the \textsc{Helmholtz} equation in particular serving as a readable alternative.%
\footnote{%
  \cite{gumerov2004fast} \textsc{Gumerov} \& \textsc{Duraiswami}, \S3
}
We shall however address first the matter of the result of the preceding text.
We shall insert into equation \eqref{eq:hypersphere_green} the relation \eqref{eq:gegenbauer-legendre} for the threedimensional case, yielding

\begin{equation}
  \green(\xvec; \bzhe) = \frac{1}{4\pi\absl{\bzhe}} \sum_{k = 0}^{\infty} \varepsilon^k \legendre_k(\cos{\alpha}).
\end{equation}

\noindent
We shall exploit the fact that the \textsc{Gegenbauer} functions evaluated at the inner product between two higherdimensional vectors can be written as a superposition of products of the spherical harmonics.%
\footnote{%
  \cite{avery2018hyperspherical} \textsc{Avery} \& \textsc{Avery}, \S3.2
}
This result simplifies further to what the literature on fast multipole calculations often labels the addition theorem for three dimensions, and is given as follows.

\begin{theorem}
  Let $\legendre_n$ denote the \textsc{Legendre} polynomial of order $n$, $\spherical_{n}^{m}$ be a spherical harmonic, and $\xvec_1$, $\xvec_2$ be points on the unit two-sphere, such that $\xvec_i = (\cos{\alpha_{i}^{1}}, \sin{\alpha_{i}^{0}} \sin{\alpha_{i}^{1}}, \cos{\alpha_{i}^{0}} \sin{\alpha_{i}^{1}})$.
  Then,
  \begin{equation}
    \legendre_n(\cos{\alpha}) = \frac{4\pi}{2n+1} \sum_{m = -n}^{n} \spherical_{n}^{m}(\alpha_{1}^{0}, \alpha_{1}^{1}) {\spherical_n^m}^{\dag}(\alpha_{2}^{0}, \alpha_{2}^{1}),
  \end{equation}
  where $\alpha$ is the angle between the points.
\end{theorem}

\noindent
We may now address the translation directly.
It shall promtply become apparent that we seek local and distal expansions of the \textsc{Green} function.
Relative to some new origin $O^{\prime}$ with vector $\ovec$ from the absolute origin, we write the point of evaluation $\bzhe^{\prime} = \bzhe - \ovec$.
The local and distal expansions shall respectively regular and singular at this new origin.
